\chapter*{Цели и задачи}
\addcontentsline{toc}{chapter}{Введение}
\label{ch:intro_lab8}

\textbf{Цель лабораторной работы:} научиться применять разработанный пайплайн для тиражирования кода с целью решения задачи полиномиальной регрессии.

\textbf{Основные задачи:}
\begin{itemize}
    \item получение навыков рефакторинга кода в проектах машинного обучения;
    \item изучение поведения модели полиномиальной регрессии при изменении степени полинома;
    \item исследование свойств набора данных в рамках задачи полиномиальной регрессии.
\end{itemize}

\textbf{Теоретическое обоснование:}

Полиномиальная регрессия является частным случаем множественной линейной регрессии, в которой связь между независимой переменной $x$ и зависимой переменной $y$ моделируется полиномом степени $n$. Модель полиномиальной регрессии имеет вид:

\begin{equation}
y = \beta_0 + \beta_1 x + \beta_2 x^2 + \ldots + \beta_n x^n + \varepsilon
\label{eq:polynomial}
\end{equation}

где $\beta_i$ — коэффициенты модели, $n$ — степень полинома, $\varepsilon$ — ошибка модели.

Преимущество полиномиальной регрессии заключается в способности аппроксимировать нелинейные зависимости между переменными. Однако при выборе степени полинома необходимо соблюдать баланс: слишком низкая степень приводит к недообучению (underfitting), а слишком высокая — к переобучению (overfitting).

В качестве набора данных для выполнения лабораторной работы используется набор данных \textbf{Position Salaries}, который содержит информацию о должностях и соответствующих уровнях и зарплатах. Данный набор демонстрирует нелинейную зависимость между уровнем должности и зарплатой.

\endinput